\section{The derived return and the complete arithmetic source
correspondence}\label{the-derived-return-and-the-complete-arithmetic-source-correspondence}

25 September 2026. Complete derivation FSC0--FSC6.

\subsection{FSC0. Why the whole source must be
checked}\label{fsc0.-why-the-whole-source-must-be-checked}

The product pairing and the derived diagonal return have been
constructed on the actual Connes--Consani receiving topology in
\href{https://github.com/KokunoYumeto/zeta-function-research-reader/blob/main/workbenches/splitzero-tandem/continuations/20260925-full-source-tensor-and-original-divisor/counterfactual/ORIGINAL_RESIDUE_PRODUCT_SHEAF_PAIRING.md}{PRS},
\href{https://github.com/KokunoYumeto/zeta-function-research-reader/blob/main/workbenches/splitzero-tandem/continuations/20260925-full-source-tensor-and-original-divisor/counterfactual/RESIDUE_PRODUCT_DIAGONAL_INDEPENDENT.md}{RPD}
and
\href{https://github.com/KokunoYumeto/zeta-function-research-reader/blob/main/workbenches/splitzero-tandem/continuations/20260925-full-source-tensor-and-original-divisor/counterfactual/DIAGONAL_EXTRAORDINARY_RETURN.md}{the
derived return}. The user's complete-spectrum argument requires a
further check: retain every original arithmetic point when returning
from that receiving base. The calculation here performs this check,
identifies the exact limit of a single-map return, and constructs the
full correspondence with its coefficient maps.

The source definitions remain
\href{https://github.com/KokunoYumeto/zeta-function-research-reader/blob/36a82ecca98addb8a98b48204e349737856c7223/workbenches/splitzero-tandem/continuations/20260924-cohomology-weight-and-character-lifts/original-character-lifting/foundations/user_definitions_20260924/SOURCE_OPERATIONS_AND_PROOFS.md}{B1--B5
and retraction R1}. The two-chart arithmetic space used here is the
specified source-induced construction
\href{https://github.com/KokunoYumeto/zeta-function-research-reader/blob/36a82ecca98addb8a98b48204e349737856c7223/workbenches/splitzero-tandem/continuations/20260924-cohomology-weight-and-character-lifts/original-character-lifting/tau_lifting_weights_20260924/SOURCE_CC_DOUBLE_PULLBACK.md}{DCP1--DCP2};
it is not itself a numerical definition of primitive \(Z_1/\tau\). All
integer and prime data used below are in the already supplied arithmetic
layer. No counting operation on \(\tau\), source addition, new parity,
or source weight is used. The user's whole-spectrum, global-quotient and
operation-prerequisite arguments are retained in the complete corpus, in
particular USR-9ad1c0a2d09dba92, USR-6152e3bc6302258c and
USR-4322be19bff532cd.

Human provenance: Alain Connes and Caterina Consani,
\href{https://arxiv.org/abs/0903.2024v3}{\emph{Schemes over
\(\mathbb F_1\) and zeta functions}, 0903.2024v3, §5}, for the receiving
base and coefficient geometry. DCP constructs the actual source
comparison used here. The full source-return calculation below is a
programme derivation, not a claim found verbatim in that paper. The
duality target remains Pierre Deligne's
\href{https://www.numdam.org/item/PMIHES_1980__52__137_0/}{\emph{La
conjecture de Weil. II}, §§3.3.11 and 3.6}, without importing its
finite-field weight assertions.

\subsection{FSC1. The exact spaces and the specified
return}\label{fsc1.-the-exact-spaces-and-the-specified-return}

Put \[
X=U\cup\{\mathfrak m_+,\mathfrak m_-\},\qquad U=\operatorname{Spec}\mathbb Z.
\] The opens of \(X\) are the ordinary opens contained in \(U\), the two
full charts \(X_\pm=U\cup\{\mathfrak m_\pm\}\), and \(X\) itself. This
is the topology already constructed in DCP; it retains each arithmetic
point separately. Let \[
f:X\to Y=\{c_+,c_-,\eta\},\qquad f(U)=\{\eta\},\quad f(\mathfrak m_\pm)=c_\pm.
\tag{FSC1.1}
\] The receiving map in the derived return is \[
q:Y^2\to Y,\qquad
q(c_+,c_+)=c_+,\quad q(c_-,c_-)=c_-,\quad
q(x,y)=\eta\text{ at the other seven points}.
\tag{FSC1.2}
\] It is continuous: the inverse image of \(U_+\) is the complement of
the closed point \((c_-,c_-)\), the inverse image of \(U_-\) is the
complement of \((c_+,c_+)\), and the inverse image of \(\{\eta\}\) is
the complement of both. Its diagonal section is \(\Delta_Y(y)=(y,y)\).

All squares and products in this note are the indicated topological
products and their sheaves. They are not identified with an arithmetic
scheme fibre product over an unstated base.

\subsection{FSC2. A return map preserving the full arithmetic diagonal
cannot lift this
q}\label{fsc2.-a-return-map-preserving-the-full-arithmetic-diagonal-cannot-lift-this-q}

There is no continuous map \(\widetilde q:X^2\to X\) satisfying both \[
f\widetilde q=q(f\times f),\qquad \widetilde q\Delta_X=\operatorname{id}_X,
\quad\Delta_X(x)=(x,x).
\tag{FSC2.1}
\] This is a precise statement about the displayed map \(q\), spaces and
two equations, not a nonexistence assertion about the user's source or
every geometric return.

To prove it, first calculate the closure of an arithmetic closed point
\(p\): \[
\overline{\{p\}}^{\,X}=\{p,\mathfrak m_+,\mathfrak m_-\}.
\tag{FSC2.2}
\] Inside \(U\), the ordinary prime point is closed. Every neighbourhood
of either \(\mathfrak m_\pm\) contains the entire \(U\), hence contains
\(p\). No other point lies in the closure, by these same open sets and
the ordinary closure in \(U\). In the product topology, the closure of a
singleton pair is the product of the two singleton closures: a basic
product neighbourhood meets that singleton exactly when its two
component neighbourhoods do. Therefore \[
(\mathfrak m_+,\mathfrak m_-)\in\overline{\{(p,p)\}}^{\,X^2}
\quad\text{for every arithmetic prime }p.
\] Continuity and the second equation of (FSC2.1) force \[
\widetilde q(\mathfrak m_+,\mathfrak m_-)
\in\bigcap_p\overline{\{p\}}^{\,X}
=\{\mathfrak m_+,\mathfrak m_-\}.
\tag{FSC2.3}
\] The last equality already follows from two distinct prime points, and
remains the intersection for the entire spectrum. But the first equation
of (FSC2.1) gives \[
f\widetilde q(\mathfrak m_+,\mathfrak m_-)=q(c_+,c_-)=\eta,
\] so that same image must belong to \(U=f^{-1}\{\eta\}\). The two
specified subsets of \(X\) are disjoint. This proves the claim from the
exact topology and map requirements.

The operation prerequisites explain its scope. The diagonal identity was
required to keep each recovered arithmetic point, including all prime
points, rather than select a constant arithmetic output. The
contradiction does not arise from adding primitive \(\tau\) or assigning
it a coordinate. It rules out only a continuous single-valued lift of
this particular receiving return with that diagonal requirement.

\subsection{FSC3. Construct the full return as a
correspondence}\label{fsc3.-construct-the-full-return-as-a-correspondence}

Define the topological fibre product \[
\mathcal Z=\{(x,y,z)\in X^2\times X:f(z)=q(f(x),f(y))\}
\tag{FSC3.1}
\] with its subspace topology. It has continuous projections \[
a:\mathcal Z\to X^2,\quad a(x,y,z)=(x,y),\qquad
b:\mathcal Z\to X,\quad b(x,y,z)=z.
\tag{FSC3.2}
\] Their exact relation is \[
fb=q(f\times f)a.
\tag{FSC3.3}
\] Every fibre of \(a\) is explicitly determined: \[
a^{-1}(\mathfrak m_+,\mathfrak m_+)=\{(\mathfrak m_+,\mathfrak m_+,\mathfrak m_+)\},
\] \[
a^{-1}(\mathfrak m_-,\mathfrak m_-)=\{(\mathfrak m_-,\mathfrak m_-,\mathfrak m_-)\},
\] \[
a^{-1}(x,y)=\{(x,y,z):z\in U\}\cong U
\quad\text{for every other pair }(x,y).
\tag{FSC3.4}
\] These are homeomorphisms onto the indicated subspace fibres, since
the first two coordinates are fixed. In particular \(a\) is onto and
every arithmetic output point is retained in every nonexceptional fibre;
the original primes have not been replaced by their image \(\eta\).

There is a continuous lifted diagonal \[
\delta_X:X\to\mathcal Z,\qquad x\mapsto(x,x,x).
\tag{FSC3.5}
\] Its image lies in \(\mathcal Z\) because
\(q\Delta_Y=\operatorname{id}_Y\). Continuity is that of the ordinary
triple diagonal followed by restriction of codomain. The identities \[
a\delta_X=\Delta_X,\qquad b\delta_X=\operatorname{id}_X
\tag{FSC3.6}
\] are pointwise. Thus the correspondence preserves every point of the
full arithmetic diagonal, exactly where a single-valued return in
(FSC2.1) fails. No closedness, properness, or scheme structure is
asserted for \(\mathcal Z\) without proof.

\subsection{FSC4. Exact coefficient morphisms on the
correspondence}\label{fsc4.-exact-coefficient-morphisms-on-the-correspondence}

For every receiving sheaf \(\mathcal G\) on \(Y\), the equality of
continuous maps (FSC3.3) induces the canonical isomorphism \[
a^{-1}(f\times f)^{-1}q^{-1}\mathcal G
\xrightarrow{\ \cong\ }b^{-1}f^{-1}\mathcal G.
\tag{FSC4.1}
\] Here is its explicit construction. At \((x,y,z)\), the left stalk is
\(\mathcal G_{q(f(x),f(y))}\) and the right stalk is
\(\mathcal G_{f(z)}\). The defining equality in (FSC3.1) identifies
these as the same stalk; the map is its identity. All restriction maps
are induced by the same composite continuous map in (FSC3.3). The
stalkwise identity maps therefore define mutually inverse sheaf
morphisms. This is also the usual composition isomorphism of
inverse-image functors, now with every map specified.

The derived-return calculation proves
\(q^{-1}\mathcal G=\Delta_{Y*}\mathcal G\), including all receiving
restrictions. Substitution in (FSC4.1) proves the exact return of that
diagonal direct image through the correspondence: \[
a^{-1}(f\times f)^{-1}\Delta_{Y*}\mathcal G
\cong b^{-1}f^{-1}\mathcal G.
\tag{FSC4.2}
\] Likewise the original product morphism
\(\beta:\mathcal F\boxtimes\mathcal F\to j_{\eta\eta!}\chi_{\rm dil}\mathbb C\)
pulls back along \((f\times f)a\) to a sheaf morphism on \(\mathcal Z\).
Its nonzero-amplitude component is exactly the original contour
(PRS1.4); neither its denominator nor any support label changes. Pulling
further along \(\delta_X\) gives precisely the ordinary diagonal
pullback already calculated, since \((f\times f)a\delta_X=\Delta_Yf\).

These identities concern the constructed inverse-image maps. They do not
assume a proper-base-change identity for the projections \(a,b\), nor
replace \(R\Delta_Y^!\) by ordinary inverse image. They construct an
explicit correspondence preserving the complete arithmetic diagonal and
every nonexceptional arithmetic fibre. FSC6 computes its derived
coefficient return directly.

\subsection{FSC5. Result and the next receiving
object}\label{fsc5.-result-and-the-next-receiving-object}

The receiver's derived diagonal return is valid and restores its
original global trace. FSC2 proves exactly why that receiver map must
not be called a continuous single-valued return of the entire arithmetic
source satisfying (FSC2.1). FSC3--FSC4 immediately constructs the full
alternative: the topological fibre product, both projections, the lifted
full diagonal, every fibre and the actual coefficient isomorphisms. This
applies the user's whole-spectrum requirement instead of discarding it
at the finite receiving base.

The resulting object for source-level duality is \(\mathcal Z\) with
\((a,b,\delta_X)\) and the coefficient morphism (FSC4.2). The next
calculation is its actual derived direct-image return. It is carried out
in FSC6 below, retaining the full arithmetic fibres. This result asserts
neither a numerical weight on primitive \(Z_1/\tau\) nor a proof or
disproof of RH.

\subsection{FSC6. The full correspondence has an exact derived
coefficient
return}\label{fsc6.-the-full-correspondence-has-an-exact-derived-coefficient-return}

Write \(r=f\times f\), \(v=qr:X^2\to Y\), and
\(g=fb=va:\mathcal Z\to Y\). For every sheaf \(\mathcal G\) of complex
vector spaces on \(Y\), the canonical unit gives an actual
quasi-isomorphism \[
\boxed{v^{-1}\mathcal G\longrightarrow Ra_*g^{-1}\mathcal G
=Ra_*b^{-1}f^{-1}\mathcal G.}
\tag{FSC6.1}
\] The same statement holds for bounded sheaf complexes, with their full
differentials. This is a derived calculation on the specified
correspondence, not an assumed general base-change theorem.

To prove it, let \(o=(0)\in\operatorname{Spec}\mathbb Z\) be the
original generic arithmetic point. Every nonempty open of \(X\) contains
\(o\). Thus every nonempty open of \(X^2\) contains \((o,o)\). The
triple \((o,o,o)\) belongs to \(\mathcal Z\), and every nonempty
relative open of \(\mathcal Z\) contains it: that open is an
intersection with an ambient open in \(X^3\), all of whose nonempty
product neighbourhoods contain the generic triple. Consequently the
constant sheaf with value any complex vector space \(V\) is exactly the
direct image of \(V\) from that generic point, on each of \(X^2\) and
\(\mathcal Z\). This assertion concerns the open-set sections, all equal
to \(V\) on nonempty opens with identity restrictions.

Such a sheaf is injective in the algebraic category of sheaves of
complex vector spaces. In fact, for any point inclusion \(i_x\), \[
\operatorname{Hom}(\mathcal H,i_{x*}V)=\operatorname{Hom}(\mathcal H_x,V).
\] The stalk functor is exact, and linear maps from a subspace into
\(V\) extend by a vector-space complement. Hence the right-hand functor
is exact in \(\mathcal H\). This proves injectivity both for the stated
generic-point constant sheaves and for the closed-point skyscrapers
below; it makes no assertion about injectivity in a category requiring
continuous linear maps.

Use the three receiving injectives \[
I_\eta(V)=(V\xrightarrow{1}V\xleftarrow{1}V),\qquad
I_+(V)=(V\to0\leftarrow0),\qquad I_-(V)=(0\to0\leftarrow V),
\tag{FSC6.2}
\] where the entries in the last two diagrams are the plus, generic and
minus entries in that order. The inverse images of \(I_\eta(V)\) by
\(v\) and \(g\) are the constant sheaves just computed. The inverse
image of \(I_+(V)\) is the skyscraper at
\((\mathfrak m_+,\mathfrak m_+)\) on \(X^2\), respectively at
\((\mathfrak m_+,\mathfrak m_+,\mathfrak m_+)\) on \(\mathcal Z\). The
minus case has the corresponding minus points. To prove these sheaf
identifications, their inverse-image stalks are \(V\) exactly at the
indicated fibre points and zero elsewhere. The adjunction map to the
skyscraper is identity at that point and an isomorphism on every other
zero stalk, and is therefore a sheaf isomorphism.

All of these inverse-image sheaves on \(\mathcal Z\) are injective by
the preceding point-inclusion argument. Moreover their direct images
along \(a\) are exactly their counterparts on \(X^2\), with the
canonical unit equal to identity. For the constant sheaf, \(a^{-1}O\) is
nonempty for each nonempty open \(O\subset X^2\) because \(a\) is onto,
and it contains the same generic triple. Its sections are \(V\), giving
the identity on every nonempty open. For a skyscraper this follows from
the literal equality of composites from the indicated point through
\(a\). There are no higher derived direct images, since each sheaf on
\(\mathcal Z\) is injective.

For completeness every receiving sheaf
\(\mathcal G=(V_+\xrightarrow{s_+}B\xleftarrow{s_-}V_-)\) has the
following finite, functorial injective resolution: \[
0\to\mathcal G\longrightarrow
I_\eta(B)\oplus I_+(V_+)\oplus I_-(V_-)
\longrightarrow I_+(B)\oplus I_-(B)\to0.
\tag{FSC6.3}
\] At the generic stalk the first map is identity and the last target is
zero. At the plus stalk the first map is \(x\mapsto(s_+x,x)\), and the
differential is \((b,x)\mapsto b-s_+x\); at the minus stalk it is
\(x\mapsto(s_-x,x)\) with differential \((b,x)\mapsto b-s_-x\). These
maps commute with the original restrictions. Their kernels and
surjectivity prove exactness at all three stalks. Every term is
injective by the same evaluation adjunction used in FTD.

Inverse image is exact for these sheaves and sends every term of
(FSC6.3) to one of the injectives already computed on \(\mathcal Z\). It
is therefore an injective resolution of \(g^{-1}\mathcal G\). Applying
\(a_*\) gives term by term the identical resolution obtained by
\(v^{-1}\) from (FSC6.3), with the same differentials and unit. This
proves (FSC6.1). For a bounded complex, use the functorial resolution
(FSC6.3) in each degree and its total complex; it is a bounded complex
of injectives after \(g^{-1}\), so the same termwise proof applies
without dropping any degree.

Combining (FSC6.1) with the exact receiving identity
\(q^{-1}=\Delta_{Y*}\) proves \[
\boxed{Ra_*b^{-1}f^{-1}\mathcal G\cong
(f\times f)^{-1}\Delta_{Y*}\mathcal G.}
\tag{FSC6.4}
\] In particular one may substitute the already constructed full complex
\(\mathcal G=Rq_*K\) or
\(\mathcal G=Rq_*(\mathcal F\boxtimes\mathcal F)\), and the actual
morphism between them. All are bounded by DER. The isomorphism is
natural in \(\mathcal G\), so it retains that morphism, every term and
the counit; it does not assert a new scalar equality merely because
cohomology dimensions agree. The full arithmetic source return is
therefore constructed as a correspondence with an exact derived
coefficient comparison. A numerical weight separation is not inserted
into that comparison.

\subsubsection{The degree-one mixed obstruction has an explicit global
equivariant
contraction}\label{the-degree-one-mixed-obstruction-has-an-explicit-global-equivariant-contraction}

GLOBAL\_MIXED\_RETURN\_CONTRACTION.md, GMC0--GMC7, follows the generic
mixed row into its complete global sheaf complex. The actual section
s\_+(j)=(Sigma\^{}\{-1\}j,0,0) is continuous and dilation equivariant,
using the already proved summation inverse. It gives
S(m\_1,m\_2)=(-(s\_+ tensor1)m\_1,(1 tensor s\_+)m\_2,0,0), with d\_M
S=id. The exact global cochain map is (id-Sd\_M,vartheta), onto K\_0 in
degree0 and Q tensor Q in degree1, with zero differential. Its entire
kernel is {[}S(M)→iota(M){]}, contracted by h(iota m)=Sm. Every original
prime action intertwines.

Thus the degree-one mixed obstruction is killed by the constructed
contraction of ker(F)={[}S(M)→iota(M){]}. The full mixed subsheaf
retains RΓ(Y,N)≃K\_0{[}0{]}, with K\_0=(H tensor Q) direct-sum (Q tensor
H); it is not asserted to be acyclic. The generic polynomial
representative P(L\_total)z=iota m becomes the exact global boundary
d(Sm). All original endpoint and extra copies remain in K\_0=(H tensor
Q) direct-sum (Q tensor H); the original residue trace factors through
the computed map with coefficient +1. This is a theorem in the specified
algebraic tensor/sheaf calculation, with no assumption of
completed-tensor cohomology, mirror equivariance of the plus-only
section, or numerical purity. The remaining Q operators and all primary
nilpotents are retained explicitly in GMC6.3.

\subsection{Publication source and dependency
links}\label{publication-source-and-dependency-links}

The source geometry is Alain Connes and Caterina Consani, \emph{Schemes
over F1 and zeta functions},
\href{https://arxiv.org/abs/0903.2024v3}{arXiv:0903.2024v3, §5}. The
weight-control comparison is Pierre Deligne, \emph{La conjecture de
Weil. II}, Publications mathematiques de l'IHES52 (1980),137-252,
\href{https://www.numdam.org/item/PMIHES_1980__52__137_0/}{§§3.3.11
and3.6}. Deligne material was read through the identified French
transcription and separately recorded peer source-page checks, not
Deligne-authored TeX. These sources are not asserted to contain the new
programme derivations. The
\href{https://github.com/KokunoYumeto/zeta-function-research-reader/blob/36a82ecca98addb8a98b48204e349737856c7223/workbenches/splitzero-tandem/continuations/20260924-cohomology-weight-and-character-lifts/BUILD_AND_REVIEW.md}{preceding
complete proof and reading record} records inherited source versions and
actual inspection limits. The investigator's corrected construction
remains attributed in the original text above.
